\section{Joint embedding prediction and the collapse question}
\label{sec:C05:jepa}
I-JEPA predicts representations of target image regions from a context representation, using a predictor conditioned on target positions and a slowly updated target encoder~\cite{assran2023ijepa}. In schematic notation, a context encoder gives $z_C=f_\phi(x_C)$, a target encoder gives $z_B=\overline f_{\bar\phi}(x)_B$, and a predictor produces $\widehat z_B=q_\psi(z_C,B)$. For illustration, a squared representation loss takes the form
\begin{equation}
 \mathcal L_{\mathrm{pred}}=
 \mathbb E_{x,C,B}\|\widehat z_B-
 \operatorname{sg}(z_B)\|_2^2.
\end{equation}
Here $\operatorname{sg}$ stops gradients through the target, and $\bar\phi$ is updated from the online encoder through an exponential moving average. The mask distribution determines which contextual and spatial relations must be learned. This squared loss is a schematic analysis objective: the inspected I-JEPA implementation uses smooth $L_1$, so reproducing that recipe requires its actual loss rather than substituting this expression.

Prediction in a learned latent space can avoid expending all capacity on unpredictable pixel detail. That is an attractive objective for agents whose actions depend on stable structure rather than every texture variation. However, which variation becomes irrelevant is learned from the data and objective. If a tiny visual indicator is weakly predictable yet action-critical, an abstraction optimized for broad semantic transfer may discard it. The appropriate question is whether the representation preserves task-relevant conditional distinctions, not whether it is more abstract in an unspecified sense.

The constant-representation solution clarifies a logical limitation. If both encoders output the same constant and the predictor outputs it, the schematic prediction loss is zero. Stop-gradient and a moving target change optimization dynamics but do not, by themselves, make this algebraic solution impossible. Actual JEPA training uses architectural, masking, normalization, and optimization choices that empirically avoid uninformative solutions in the reported settings. The existence of useful trained representations is an empirical result; a universal noncollapse guarantee should not be inferred from the abbreviated objective. This distinction is especially important when adapting the method to a new modality or action-conditioned domain.

V-JEPA extends feature prediction to video representations, learning temporal structure without reconstructing every observed pixel~\cite{bardes2024vjepa}. Temporal predictability and interventional identification remain different properties. Consider a binary latent variable $U\sim\operatorname{Bernoulli}(1/2)$, an observational action $A=U$, and two possible transition mechanisms: $Y=A$ or $Y=U$. Both generate exactly the same observational pairs $(A,Y)$. Under $\operatorname{do}(A=1)$, however, the first gives $P(Y=1)=1$ and the second gives $P(Y=1)=1/2$. This constructed counterexample shows that perfect passive prediction does not identify the effect of an intervention absent additional assumptions or action-grounded evidence.

The 2026 progression from dense V-JEPA features to action-conditioned policy components should therefore be evaluated at two distinct interfaces: preservation of observation distinctions and prediction of consequences under the deployed action distribution~\cite{rewrite_vjepa21,rewrite_vjepapolicy2026}. A changed action schema can break the second while leaving the first intact. Relevant controls retain the encoder while changing the policy or action expert, then retain the action interface while changing the encoder. A frozen visual backbone can supply useful state features without constituting an identified transition model or a complete controller.

